\begin{figure}[H]
\centering
\begin{tikzpicture}[>=Stealth,scale=0.88]
  \coordinate (Obj) at (0,0);
  \coordinate (Foc) at (4.6,0);
  \coordinate (Ocu) at (6.3,0);

  % Eje óptico.
  \draw[->,gray!75] (-5.1,0) -- (9,0) node[right,black] {eje óptico};

  % Objetivo convergente.
  \filldraw[fill=cyan!18,draw=blue!65!black,thick]
    (0,-2.15) .. controls (-0.68,-1.05) and (-0.68,1.05) .. (0,2.15)
    .. controls (0.68,1.05) and (0.68,-1.05) .. cycle;
  \node[above] at (0,2.3) {objetivo};
  \draw[<->,thick] (-0.88,-2.15) -- (-0.88,2.15)
    node[midway,left,fill=white,inner sep=1.5pt] {$D$};

  % Rayos paralelos procedentes de un objeto lejano.
  \foreach \y in {-1.25,0,1.25}
    \draw[->,blue,thick] (-4.8,\y) -- (0,\y);
  \node[blue!70!black,align=center] at (-3.25,1.72)
    {rayos paralelos\\del objeto lejano};

  % El objetivo forma la imagen primaria en su plano focal.
  \draw[red!80!black,thick] (0,1.25) -- (Foc) -- (6.3,-0.462);
  \draw[red!80!black,thick] (0,0) -- (Foc) -- (Ocu);
  \draw[red!80!black,thick] (0,-1.25) -- (Foc) -- (6.3,0.462);
  \draw[dashed,gray!75] (4.6,-1.55) -- (4.6,1.55);
  \fill[red!75!black] (Foc) circle (2.2pt);
  \draw[->,very thick,orange!85!black] (4.6,0.72) -- (4.6,0);
  \node[align=center,above] at (4.6,1.62)
    {plano focal\\imagen primaria};

  % Ocular: la imagen primaria está situada en su plano focal anterior.
  \filldraw[fill=cyan!15,draw=cyan!45!blue,thick]
    (6.3,-1.25) .. controls (5.92,-0.62) and (5.92,0.62) .. (6.3,1.25)
    .. controls (6.68,0.62) and (6.68,-0.62) .. cycle;
  \node[above] at (6.3,1.4) {ocular};

  % Para observación relajada, el ocular entrega un haz aproximadamente colimado.
  \draw[->,green!55!black,thick] (6.3,0.462) -- (8.15,0.462);
  \draw[->,green!55!black,thick] (6.3,0) -- (8.15,0);
  \draw[->,green!55!black,thick] (6.3,-0.462) -- (8.15,-0.462);
  \node[green!45!black,above] at (7.3,0.55) {haz colimado};

  % Ojo del observador.
  \draw[thick] (8.45,0) ellipse (0.3 and 0.66);
  \fill (8.17,0) circle (0.11);
  \node[below] at (8.45,-0.7) {ojo};

  % Distancias focales.
  \draw[gray!55] (0,-2.25) -- (0,-2.95);
  \draw[gray!55] (4.6,-1.65) -- (4.6,-2.95);
  \draw[gray!55] (6.3,-1.35) -- (6.3,-2.45);
  \draw[<->] (0,-2.82) -- (4.6,-2.82)
    node[midway,fill=white,inner sep=2pt] {$f_{\mathrm{obj}}$};
  \draw[<->] (4.6,-2.25) -- (6.3,-2.25)
    node[midway,fill=white,inner sep=2pt] {$f_{\mathrm{oc}}$};

\end{tikzpicture}
\caption{Esquema básico de un telescopio refractor ajustado para observación relajada: el objetivo forma la imagen primaria y el ocular entrega un haz aproximadamente colimado al ojo.}
\end{figure}